\setcounter{chapter}{-1}
\chapter{Зачем это RIN: персонаж не как маска, а как интерфейс личности}

\setcounter{section}{0}
\section{Почему "аниме-архетип" может быть сильнее, чем попытка реалистичного 3D-компаньона}

Есть соблазнительная интуиция, что чем ближе цифровой компаньон к человеку, тем он живее: больше полигонов, точнее мимика, реалистичнее голос, движения без сглаживания. Кажется, что реализм — прямая дорога к правдоподобию. На практике происходит обратное, и об этом стоит знать заранее, а не набивать шишку самим: чем ближе виртуальное лицо подбирается к настоящему, не дотягивая до него окончательно, тем сильнее реакция становится не эмпатией, а тревогой. Ещё в 1970 году Масахиро Мори дал этому провалу доверия имя — uncanny valley — и с тех пор оно только подтверждается: ровно там, где сходство почти полное, но не полное, не-совсем-живые глаза и не-совсем-естественная улыбка читаются не как "почти человек", а как "что-то не так с этим человеком".

Стоит разобрать не только сам эффект, но и почему он вообще возникает, а не просто принять его как данность. Человеческое восприятие лица — не нейтральный сканер форм, а система, натренированная (той же самой способностью, о которой в 1.1 пойдёт речь применительно к маскам но) обнаруживать мельчайшие отклонения от знакомого паттерна, потому что именно мельчайшее отклонение в лице другого человека веками сигнализировало об угрозе, болезни, обмане. Одно из правдоподобных объяснений в том, что эта система, заточенная ловить подделку, встречая почти точную подделку, не расслабляется, а включается на полную мощность: почти-точность читается ею как искусно скрытая угроза, а не как её отсутствие. Обидная деталь в том, что реализм не обходит эту яму стороной, а в неё и ведёт, потому что чем детальнее попытка, тем больше у этой системы материала для срабатывания.

Аниме-архетип решает проблему не тем, что маскирует зазор, а тем, что вообще не пытается его закрыть. Rei Ayanami не претендует на фотореализм — её дизайн условен, и палитра говорит раньше, чем она открывает рот: бледная кожа, синие волосы, красные глаза читаются как "нечеловеческая тишина" ещё до первой реплики. Условность не обманывает зрителя, будто перед ним живой человек, — она прямо заявляет: это персонаж, интерфейс, знак. И именно поэтому зрителю не с чем спорить на уровне "а почему у неё веки чуть-чуть неправильно двигаются". Она обходит проверку на реализм, потому что никогда её не заявляла.

Вот первый пункт на карте всей книги, и дальше мы будем находить его снова и снова в разных костюмах: ограничение средства выражения — не помеха, если зритель заранее не ждёт от этого средства того, чего оно дать не может. Для RIN это не эстетическое предпочтение, а прямое снятие риска. У проекта нет и в обозримой перспективе не будет фотореалистичного тела, точной лицевой анимации, настоящего голоса без цифровых сглаживаний. Если целиться в "почти человека", единственный доступный результат — застрять где-то в долине между живым и мёртвым. Если целиться в убедительный условный образ — ночная терминальная ведьма, характерный голос, узнаваемый жест, — зазор с реальностью перестаёт быть уязвимостью, потому что зритель никогда не ждал, что зазора не будет.
